% MechVal Scorecard — standalone \input{} file
% Requires: booktabs, array, multirow, xcolor, amssymb packages in parent document
%
% Usage: % MechVal Scorecard — standalone \input{} file
% Requires: booktabs, array, multirow, xcolor, amssymb packages in parent document
%
% Usage: % MechVal Scorecard — standalone \input{} file
% Requires: booktabs, array, multirow, xcolor, amssymb packages in parent document
%
% Usage: % MechVal Scorecard — standalone \input{} file
% Requires: booktabs, array, multirow, xcolor, amssymb packages in parent document
%
% Usage: \input{mechval_scorecard}

\definecolor{confirmed}{HTML}{2E7D32}
\definecolor{partial}{HTML}{E65100}
\definecolor{untested}{HTML}{757575}
\definecolor{inconclusive}{HTML}{9E9E9E}

\providecommand{\cmark}{{\color{confirmed}\checkmark}}%
\providecommand{\pmark}{{\color{partial}$\circ\mkern-7mu/$}}%
\providecommand{\umark}{{\color{untested}?}}%
\providecommand{\xmark}{{\color{red}$\times$}}%

\section{Mechanistic Validity Scorecard}
\label{app:mechval}

Mechanistic Validity \citep{tower2026mechval} provides a systematic framework
for evaluating whether mechanistic claims about neural networks are warranted
by available evidence. The framework organizes 30 criteria into five types:
\emph{construct} validity (whether the claimed mechanism is well-defined and
falsifiable), \emph{measurement} validity (whether the metrics reliably detect
the mechanism), \emph{internal} validity (whether interventions establish
necessity and sufficiency), \emph{external} validity (whether findings
generalize across prompts, tasks, and models), and \emph{interpretive} validity
(whether the level of description matches the evidence). Each criterion
receives one of four ratings: \cmark~Confirmed, \pmark~Partially confirmed,
\umark~Untested, or \xmark~Disconfirmed. The overall verdict reflects the
strongest tier of evidence reached: \emph{Falsifiable} (construct criteria
met), \emph{Causally Suggestive} (some internal criteria met),
\emph{Mechanistically Supported} (necessity and sufficiency established),
\emph{Triangulated} (multiple independent evidence lines converge), or
\emph{Validated} (cross-model replication with independent methods).

We evaluate both central claims of this paper. Claim~1 (the Grassmannian atlas
partitions operations into three geometric classes) reaches \textbf{Triangulated}:
equivariance testing, SVD prediction, and training dynamics provide independent
converging evidence. Claim~2 (structured pi-SAE recovers nonlinear causal
variables) reaches \textbf{Mechanistically Supported}: patching establishes
necessity and sufficiency, though independent non-IIA validation is still needed.

\vspace{1em}
\begin{table}[H]
\centering
\caption{Mechanistic Validity scorecard (Tower, 2026). Each criterion is rated as
\cmark~Confirmed, \pmark~Partially confirmed, \umark~Untested, or
\xmark~Disconfirmed. Claim~1 reaches \textbf{Triangulated} (multiple
converging independent evidence lines); Claim~2 reaches
\textbf{Mechanistically Supported} (necessity + sufficiency established).}
\label{tab:mechval}
\small
\setlength{\tabcolsep}{4pt}
\begin{tabular}{@{}llcc@{}}
\toprule
\textbf{Type} & \textbf{Criterion} & \textbf{Claim 1: Atlas} & \textbf{Claim 2: pi-SAE} \\
\midrule
\multirow{5}{*}{\rotatebox[origin=c]{90}{\textsc{Construct}}}
 & C1~~Falsifiability          & \cmark & \cmark \\
 & C2~~Structural plausibility & \cmark & \cmark \\
 & C3~~Convergent validity     & \cmark & \pmark \\
 & C4~~Discriminant validity   & \cmark & \cmark \\
 & C5~~Nomological validity    & \pmark & \pmark \\
\midrule
\multirow{6}{*}{\rotatebox[origin=c]{90}{\textsc{Measurement}}}
 & M1~~Reliability             & \pmark & \pmark \\
 & M2~~Baseline separation     & \cmark & \cmark \\
 & M3~~Stability               & \pmark & \pmark \\
 & M4~~Calibration             & \cmark & \cmark \\
 & M5~~Sensitivity             & \cmark & \cmark \\
 & M6~~Invariance              & \pmark & \pmark \\
\midrule
\multirow{8}{*}{\rotatebox[origin=c]{90}{\textsc{Internal}}}
 & I1~~Necessity               & \cmark & \cmark \\
 & I2~~Sufficiency             & \pmark & \cmark \\
 & I3~~Specificity             & \cmark & \cmark \\
 & I4~~Double dissociation     & \cmark & \pmark \\
 & I5~~Confound control        & \pmark & \cmark \\
 & I6~~Epistatic interaction   & \pmark & \cmark \\
 & I7~~Rescue reversibility    & \pmark & \umark \\
 & I10~Confounding sensitivity & \pmark & \pmark \\
\midrule
\multirow{6}{*}{\rotatebox[origin=c]{90}{\textsc{External}}}
 & E1~~Intervention reach      & \cmark & \cmark \\
 & E2~~Prompt generalization   & \pmark & \cmark \\
 & E3~~Cross-task               & \cmark & \cmark \\
 & E4~~Cross-model             & \umark & \pmark \\
 & E5~~Graded response         & \cmark & \cmark \\
 & E6~~Novel prediction        & \cmark & \pmark \\
\midrule
\multirow{5}{*}{\rotatebox[origin=c]{90}{\textsc{Interp.}}}
 & V1~~Level declaration       & \cmark & \cmark \\
 & V2~~Level-evidence match    & \cmark & \cmark \\
 & V3~~Alternative level       & \pmark & \pmark \\
 & V4~~Anthropomorphism check  & \cmark & \cmark \\
 & V5~~Scope declaration       & \pmark & \pmark \\
\midrule
\multicolumn{2}{@{}l}{\textbf{Confirmed / Partial / Untested}} & 17 / 12 / 1 & 18 / 11 / 1 \\
\multicolumn{2}{@{}l}{\textbf{Verdict}} & \textbf{Triangulated} & \textbf{Mech.\ Supported} \\
\bottomrule
\end{tabular}

\vspace{0.5em}
\raggedright\footnotesize
\textbf{Key gaps.}
Claim~1: cross-model generalization untested (single architecture);
stochastic grokking rests on 2 seeds; SVD vs.\ DAS not dimension-matched.
Claim~2: no independent validation beyond IIA paradigm;
hyperparameter sensitivity not ablated; VAE training seed variance not reported.
\textbf{Advancing to Validated} requires cross-architecture atlas replication
and an independent (non-IIA) validation of the pi-SAE latent structure.
\end{table}


\definecolor{confirmed}{HTML}{2E7D32}
\definecolor{partial}{HTML}{E65100}
\definecolor{untested}{HTML}{757575}
\definecolor{inconclusive}{HTML}{9E9E9E}

\providecommand{\cmark}{{\color{confirmed}\checkmark}}%
\providecommand{\pmark}{{\color{partial}$\circ\mkern-7mu/$}}%
\providecommand{\umark}{{\color{untested}?}}%
\providecommand{\xmark}{{\color{red}$\times$}}%

\section{Mechanistic Validity Scorecard}
\label{app:mechval}

Mechanistic Validity \citep{tower2026mechval} provides a systematic framework
for evaluating whether mechanistic claims about neural networks are warranted
by available evidence. The framework organizes 30 criteria into five types:
\emph{construct} validity (whether the claimed mechanism is well-defined and
falsifiable), \emph{measurement} validity (whether the metrics reliably detect
the mechanism), \emph{internal} validity (whether interventions establish
necessity and sufficiency), \emph{external} validity (whether findings
generalize across prompts, tasks, and models), and \emph{interpretive} validity
(whether the level of description matches the evidence). Each criterion
receives one of four ratings: \cmark~Confirmed, \pmark~Partially confirmed,
\umark~Untested, or \xmark~Disconfirmed. The overall verdict reflects the
strongest tier of evidence reached: \emph{Falsifiable} (construct criteria
met), \emph{Causally Suggestive} (some internal criteria met),
\emph{Mechanistically Supported} (necessity and sufficiency established),
\emph{Triangulated} (multiple independent evidence lines converge), or
\emph{Validated} (cross-model replication with independent methods).

We evaluate both central claims of this paper. Claim~1 (the Grassmannian atlas
partitions operations into three geometric classes) reaches \textbf{Triangulated}:
equivariance testing, SVD prediction, and training dynamics provide independent
converging evidence. Claim~2 (structured pi-SAE recovers nonlinear causal
variables) reaches \textbf{Mechanistically Supported}: patching establishes
necessity and sufficiency, though independent non-IIA validation is still needed.

\vspace{1em}
\begin{table}[H]
\centering
\caption{Mechanistic Validity scorecard (Tower, 2026). Each criterion is rated as
\cmark~Confirmed, \pmark~Partially confirmed, \umark~Untested, or
\xmark~Disconfirmed. Claim~1 reaches \textbf{Triangulated} (multiple
converging independent evidence lines); Claim~2 reaches
\textbf{Mechanistically Supported} (necessity + sufficiency established).}
\label{tab:mechval}
\small
\setlength{\tabcolsep}{4pt}
\begin{tabular}{@{}llcc@{}}
\toprule
\textbf{Type} & \textbf{Criterion} & \textbf{Claim 1: Atlas} & \textbf{Claim 2: pi-SAE} \\
\midrule
\multirow{5}{*}{\rotatebox[origin=c]{90}{\textsc{Construct}}}
 & C1~~Falsifiability          & \cmark & \cmark \\
 & C2~~Structural plausibility & \cmark & \cmark \\
 & C3~~Convergent validity     & \cmark & \pmark \\
 & C4~~Discriminant validity   & \cmark & \cmark \\
 & C5~~Nomological validity    & \pmark & \pmark \\
\midrule
\multirow{6}{*}{\rotatebox[origin=c]{90}{\textsc{Measurement}}}
 & M1~~Reliability             & \pmark & \pmark \\
 & M2~~Baseline separation     & \cmark & \cmark \\
 & M3~~Stability               & \pmark & \pmark \\
 & M4~~Calibration             & \cmark & \cmark \\
 & M5~~Sensitivity             & \cmark & \cmark \\
 & M6~~Invariance              & \pmark & \pmark \\
\midrule
\multirow{8}{*}{\rotatebox[origin=c]{90}{\textsc{Internal}}}
 & I1~~Necessity               & \cmark & \cmark \\
 & I2~~Sufficiency             & \pmark & \cmark \\
 & I3~~Specificity             & \cmark & \cmark \\
 & I4~~Double dissociation     & \cmark & \pmark \\
 & I5~~Confound control        & \pmark & \cmark \\
 & I6~~Epistatic interaction   & \pmark & \cmark \\
 & I7~~Rescue reversibility    & \pmark & \umark \\
 & I10~Confounding sensitivity & \pmark & \pmark \\
\midrule
\multirow{6}{*}{\rotatebox[origin=c]{90}{\textsc{External}}}
 & E1~~Intervention reach      & \cmark & \cmark \\
 & E2~~Prompt generalization   & \pmark & \cmark \\
 & E3~~Cross-task               & \cmark & \cmark \\
 & E4~~Cross-model             & \umark & \pmark \\
 & E5~~Graded response         & \cmark & \cmark \\
 & E6~~Novel prediction        & \cmark & \pmark \\
\midrule
\multirow{5}{*}{\rotatebox[origin=c]{90}{\textsc{Interp.}}}
 & V1~~Level declaration       & \cmark & \cmark \\
 & V2~~Level-evidence match    & \cmark & \cmark \\
 & V3~~Alternative level       & \pmark & \pmark \\
 & V4~~Anthropomorphism check  & \cmark & \cmark \\
 & V5~~Scope declaration       & \pmark & \pmark \\
\midrule
\multicolumn{2}{@{}l}{\textbf{Confirmed / Partial / Untested}} & 17 / 12 / 1 & 18 / 11 / 1 \\
\multicolumn{2}{@{}l}{\textbf{Verdict}} & \textbf{Triangulated} & \textbf{Mech.\ Supported} \\
\bottomrule
\end{tabular}

\vspace{0.5em}
\raggedright\footnotesize
\textbf{Key gaps.}
Claim~1: cross-model generalization untested (single architecture);
stochastic grokking rests on 2 seeds; SVD vs.\ DAS not dimension-matched.
Claim~2: no independent validation beyond IIA paradigm;
hyperparameter sensitivity not ablated; VAE training seed variance not reported.
\textbf{Advancing to Validated} requires cross-architecture atlas replication
and an independent (non-IIA) validation of the pi-SAE latent structure.
\end{table}


\definecolor{confirmed}{HTML}{2E7D32}
\definecolor{partial}{HTML}{E65100}
\definecolor{untested}{HTML}{757575}
\definecolor{inconclusive}{HTML}{9E9E9E}

\providecommand{\cmark}{{\color{confirmed}\checkmark}}%
\providecommand{\pmark}{{\color{partial}$\circ\mkern-7mu/$}}%
\providecommand{\umark}{{\color{untested}?}}%
\providecommand{\xmark}{{\color{red}$\times$}}%

\section{Mechanistic Validity Scorecard}
\label{app:mechval}

Mechanistic Validity \citep{tower2026mechval} provides a systematic framework
for evaluating whether mechanistic claims about neural networks are warranted
by available evidence. The framework organizes 30 criteria into five types:
\emph{construct} validity (whether the claimed mechanism is well-defined and
falsifiable), \emph{measurement} validity (whether the metrics reliably detect
the mechanism), \emph{internal} validity (whether interventions establish
necessity and sufficiency), \emph{external} validity (whether findings
generalize across prompts, tasks, and models), and \emph{interpretive} validity
(whether the level of description matches the evidence). Each criterion
receives one of four ratings: \cmark~Confirmed, \pmark~Partially confirmed,
\umark~Untested, or \xmark~Disconfirmed. The overall verdict reflects the
strongest tier of evidence reached: \emph{Falsifiable} (construct criteria
met), \emph{Causally Suggestive} (some internal criteria met),
\emph{Mechanistically Supported} (necessity and sufficiency established),
\emph{Triangulated} (multiple independent evidence lines converge), or
\emph{Validated} (cross-model replication with independent methods).

We evaluate both central claims of this paper. Claim~1 (the Grassmannian atlas
partitions operations into three geometric classes) reaches \textbf{Triangulated}:
equivariance testing, SVD prediction, and training dynamics provide independent
converging evidence. Claim~2 (structured pi-SAE recovers nonlinear causal
variables) reaches \textbf{Mechanistically Supported}: patching establishes
necessity and sufficiency, though independent non-IIA validation is still needed.

\vspace{1em}
\begin{table}[H]
\centering
\caption{Mechanistic Validity scorecard (Tower, 2026). Each criterion is rated as
\cmark~Confirmed, \pmark~Partially confirmed, \umark~Untested, or
\xmark~Disconfirmed. Claim~1 reaches \textbf{Triangulated} (multiple
converging independent evidence lines); Claim~2 reaches
\textbf{Mechanistically Supported} (necessity + sufficiency established).}
\label{tab:mechval}
\small
\setlength{\tabcolsep}{4pt}
\begin{tabular}{@{}llcc@{}}
\toprule
\textbf{Type} & \textbf{Criterion} & \textbf{Claim 1: Atlas} & \textbf{Claim 2: pi-SAE} \\
\midrule
\multirow{5}{*}{\rotatebox[origin=c]{90}{\textsc{Construct}}}
 & C1~~Falsifiability          & \cmark & \cmark \\
 & C2~~Structural plausibility & \cmark & \cmark \\
 & C3~~Convergent validity     & \cmark & \pmark \\
 & C4~~Discriminant validity   & \cmark & \cmark \\
 & C5~~Nomological validity    & \pmark & \pmark \\
\midrule
\multirow{6}{*}{\rotatebox[origin=c]{90}{\textsc{Measurement}}}
 & M1~~Reliability             & \pmark & \pmark \\
 & M2~~Baseline separation     & \cmark & \cmark \\
 & M3~~Stability               & \pmark & \pmark \\
 & M4~~Calibration             & \cmark & \cmark \\
 & M5~~Sensitivity             & \cmark & \cmark \\
 & M6~~Invariance              & \pmark & \pmark \\
\midrule
\multirow{8}{*}{\rotatebox[origin=c]{90}{\textsc{Internal}}}
 & I1~~Necessity               & \cmark & \cmark \\
 & I2~~Sufficiency             & \pmark & \cmark \\
 & I3~~Specificity             & \cmark & \cmark \\
 & I4~~Double dissociation     & \cmark & \pmark \\
 & I5~~Confound control        & \pmark & \cmark \\
 & I6~~Epistatic interaction   & \pmark & \cmark \\
 & I7~~Rescue reversibility    & \pmark & \umark \\
 & I10~Confounding sensitivity & \pmark & \pmark \\
\midrule
\multirow{6}{*}{\rotatebox[origin=c]{90}{\textsc{External}}}
 & E1~~Intervention reach      & \cmark & \cmark \\
 & E2~~Prompt generalization   & \pmark & \cmark \\
 & E3~~Cross-task               & \cmark & \cmark \\
 & E4~~Cross-model             & \umark & \pmark \\
 & E5~~Graded response         & \cmark & \cmark \\
 & E6~~Novel prediction        & \cmark & \pmark \\
\midrule
\multirow{5}{*}{\rotatebox[origin=c]{90}{\textsc{Interp.}}}
 & V1~~Level declaration       & \cmark & \cmark \\
 & V2~~Level-evidence match    & \cmark & \cmark \\
 & V3~~Alternative level       & \pmark & \pmark \\
 & V4~~Anthropomorphism check  & \cmark & \cmark \\
 & V5~~Scope declaration       & \pmark & \pmark \\
\midrule
\multicolumn{2}{@{}l}{\textbf{Confirmed / Partial / Untested}} & 17 / 12 / 1 & 18 / 11 / 1 \\
\multicolumn{2}{@{}l}{\textbf{Verdict}} & \textbf{Triangulated} & \textbf{Mech.\ Supported} \\
\bottomrule
\end{tabular}

\vspace{0.5em}
\raggedright\footnotesize
\textbf{Key gaps.}
Claim~1: cross-model generalization untested (single architecture);
stochastic grokking rests on 2 seeds; SVD vs.\ DAS not dimension-matched.
Claim~2: no independent validation beyond IIA paradigm;
hyperparameter sensitivity not ablated; VAE training seed variance not reported.
\textbf{Advancing to Validated} requires cross-architecture atlas replication
and an independent (non-IIA) validation of the pi-SAE latent structure.
\end{table}


\definecolor{confirmed}{HTML}{2E7D32}
\definecolor{partial}{HTML}{E65100}
\definecolor{untested}{HTML}{757575}
\definecolor{inconclusive}{HTML}{9E9E9E}

\providecommand{\cmark}{{\color{confirmed}\checkmark}}%
\providecommand{\pmark}{{\color{partial}$\circ\mkern-7mu/$}}%
\providecommand{\umark}{{\color{untested}?}}%
\providecommand{\xmark}{{\color{red}$\times$}}%

\section{Mechanistic Validity Scorecard}
\label{app:mechval}

Mechanistic Validity \citep{tower2026mechval} provides a systematic framework
for evaluating whether mechanistic claims about neural networks are warranted
by available evidence. The framework organizes 30 criteria into five types:
\emph{construct} validity (whether the claimed mechanism is well-defined and
falsifiable), \emph{measurement} validity (whether the metrics reliably detect
the mechanism), \emph{internal} validity (whether interventions establish
necessity and sufficiency), \emph{external} validity (whether findings
generalize across prompts, tasks, and models), and \emph{interpretive} validity
(whether the level of description matches the evidence). Each criterion
receives one of four ratings: \cmark~Confirmed, \pmark~Partially confirmed,
\umark~Untested, or \xmark~Disconfirmed. The overall verdict reflects the
strongest tier of evidence reached: \emph{Falsifiable} (construct criteria
met), \emph{Causally Suggestive} (some internal criteria met),
\emph{Mechanistically Supported} (necessity and sufficiency established),
\emph{Triangulated} (multiple independent evidence lines converge), or
\emph{Validated} (cross-model replication with independent methods).

We evaluate both central claims of this paper. Claim~1 (the Grassmannian atlas
partitions operations into three geometric classes) reaches \textbf{Triangulated}:
equivariance testing, SVD prediction, and training dynamics provide independent
converging evidence. Claim~2 (structured pi-SAE recovers nonlinear causal
variables) reaches \textbf{Mechanistically Supported}: patching establishes
necessity and sufficiency, though independent non-IIA validation is still needed.

\vspace{1em}
\begin{table}[H]
\centering
\caption{Mechanistic Validity scorecard (Tower, 2026). Each criterion is rated as
\cmark~Confirmed, \pmark~Partially confirmed, \umark~Untested, or
\xmark~Disconfirmed. Claim~1 reaches \textbf{Triangulated} (multiple
converging independent evidence lines); Claim~2 reaches
\textbf{Mechanistically Supported} (necessity + sufficiency established).}
\label{tab:mechval}
\small
\setlength{\tabcolsep}{4pt}
\begin{tabular}{@{}llcc@{}}
\toprule
\textbf{Type} & \textbf{Criterion} & \textbf{Claim 1: Atlas} & \textbf{Claim 2: pi-SAE} \\
\midrule
\multirow{5}{*}{\rotatebox[origin=c]{90}{\textsc{Construct}}}
 & C1~~Falsifiability          & \cmark & \cmark \\
 & C2~~Structural plausibility & \cmark & \cmark \\
 & C3~~Convergent validity     & \cmark & \pmark \\
 & C4~~Discriminant validity   & \cmark & \cmark \\
 & C5~~Nomological validity    & \pmark & \pmark \\
\midrule
\multirow{6}{*}{\rotatebox[origin=c]{90}{\textsc{Measurement}}}
 & M1~~Reliability             & \pmark & \pmark \\
 & M2~~Baseline separation     & \cmark & \cmark \\
 & M3~~Stability               & \pmark & \pmark \\
 & M4~~Calibration             & \cmark & \cmark \\
 & M5~~Sensitivity             & \cmark & \cmark \\
 & M6~~Invariance              & \pmark & \pmark \\
\midrule
\multirow{8}{*}{\rotatebox[origin=c]{90}{\textsc{Internal}}}
 & I1~~Necessity               & \cmark & \cmark \\
 & I2~~Sufficiency             & \pmark & \cmark \\
 & I3~~Specificity             & \cmark & \cmark \\
 & I4~~Double dissociation     & \cmark & \pmark \\
 & I5~~Confound control        & \pmark & \cmark \\
 & I6~~Epistatic interaction   & \pmark & \cmark \\
 & I7~~Rescue reversibility    & \pmark & \umark \\
 & I10~Confounding sensitivity & \pmark & \pmark \\
\midrule
\multirow{6}{*}{\rotatebox[origin=c]{90}{\textsc{External}}}
 & E1~~Intervention reach      & \cmark & \cmark \\
 & E2~~Prompt generalization   & \pmark & \cmark \\
 & E3~~Cross-task               & \cmark & \cmark \\
 & E4~~Cross-model             & \umark & \pmark \\
 & E5~~Graded response         & \cmark & \cmark \\
 & E6~~Novel prediction        & \cmark & \pmark \\
\midrule
\multirow{5}{*}{\rotatebox[origin=c]{90}{\textsc{Interp.}}}
 & V1~~Level declaration       & \cmark & \cmark \\
 & V2~~Level-evidence match    & \cmark & \cmark \\
 & V3~~Alternative level       & \pmark & \pmark \\
 & V4~~Anthropomorphism check  & \cmark & \cmark \\
 & V5~~Scope declaration       & \pmark & \pmark \\
\midrule
\multicolumn{2}{@{}l}{\textbf{Confirmed / Partial / Untested}} & 17 / 12 / 1 & 18 / 11 / 1 \\
\multicolumn{2}{@{}l}{\textbf{Verdict}} & \textbf{Triangulated} & \textbf{Mech.\ Supported} \\
\bottomrule
\end{tabular}

\vspace{0.5em}
\raggedright\footnotesize
\textbf{Key gaps.}
Claim~1: cross-model generalization untested (single architecture);
stochastic grokking rests on 2 seeds; SVD vs.\ DAS not dimension-matched.
Claim~2: no independent validation beyond IIA paradigm;
hyperparameter sensitivity not ablated; VAE training seed variance not reported.
\textbf{Advancing to Validated} requires cross-architecture atlas replication
and an independent (non-IIA) validation of the pi-SAE latent structure.
\end{table}
